

\textbf{Predicting unique genetic variant counts across multiple populations.}
Besides inferring trajectory by sequencing at single cell level, similar data scarcity problems arise from searching for novel genetic variants while prior knowledge of growth exists. Collecting genomics data across multiple heterogeneous populations (e.g., across different cancer types) has the potential to improve our understanding of disease. Despite sequencing advances, though, resources often remain a constraint when gathering data. So it would be useful for experimental design if experimenters with access to a pilot study of genetic sequences could predict the number of new genetic variants they might expect to find in a follow-up study: both the number of new variants shared between the populations and the total across the populations. The single-population case represents a well-studied instance of a feature allocation problem that is natural to use a flexible Bayesian nonparametric solution to incorporate prior knowledge like power-law like growth. But we here show that a natural extension of existing single-population methods can fare poorly across multiple populations that are heterogeneous, for fundamental reasons. We provide the first predictor for the number of new shared variants and new total variants that can handle heterogeneity in multiple populations. We show that our proposed method works well empirically using real cancer and population genetics data. The work \cite{shen2022double_trouble_journal} was accepted by \textit{Bayesian Analysis} with me as the sole first author. 

